\chapter{Skin Graft}

\section*{Introduction \& Clinical Indications}
A skin graft is a surgical procedure that involves the transplantation of healthy skin from a donor site to a recipient site that has experienced significant skin loss or damage. This procedure is critical in managing various clinical scenarios, including extensive burns, chronic ulcers, traumatic wounds, and defects following tumor excision. The skin graft serves to restore the protective barrier function of the skin, preventing infection, minimizing fluid loss, and facilitating both functional and aesthetic recovery of the affected area. Skin grafting is a fundamental reconstructive technique utilized across multiple surgical specialties, including plastic, general, trauma, and dermatologic surgery, to achieve durable wound closure and improve patient outcomes.

\begin{itemize}
  \item Autologous Skin Grafting
  \item Dermatoplasty
  \item Free Skin Graft
  \item Split-Thickness Skin Graft (STSG)
  \item Full-Thickness Skin Graft (FTSG)
  \item Mesh Grafting
\end{itemize}

The fundamental principle of skin grafting is the transfer of viable skin tissue from a healthy donor site to a prepared recipient wound bed, where it must establish a new blood supply to survive. This process, known as "graft take," occurs in several phases. Initially, the graft adheres to the wound bed through a fibrin layer, followed by a phase of "imbibition," where the graft passively absorbs nutrients and oxygen from the underlying wound bed. Over the next few days, "inosculation" occurs, involving the direct anastomosis of existing graft vessels with newly formed capillaries from the recipient bed, leading to active revascularization. The overarching technical goals are to achieve complete graft viability, minimize donor site morbidity, and optimize functional and aesthetic outcomes for the patient.

The choice between different types of grafts depends on the specific clinical scenario and desired outcome. Split-thickness skin grafts (STSGs) are thinner and have a higher rate of successful take, particularly on less ideal wound beds. They are harvested using a dermatome and can be meshed to expand their surface area. Full-thickness skin grafts (FTSGs), while more robust and aesthetically favorable, require a highly vascularized recipient bed and are limited by the amount of skin available for primary closure at the donor site. The surgical rationale emphasizes restoring skin integrity and function while minimizing complications.

The practice of skin grafting has ancient origins, with the earliest documented techniques found in the Sushruta Samhita, an ancient Indian text dating back to around 800 BC, which described nasal reconstruction using skin from the cheek. Modern skin grafting techniques began to take shape in the 19th century. In 1869, Swiss surgeon Jacques-Louis Reverdin introduced "pinch grafts," small epidermal grafts for covering granulating wounds. This was followed by Louis Ollier in 1872, who described larger grafts that included dermis, and Carl Thiersch, who in 1874 popularized the use of large, thin split-thickness grafts.

The early 20th century saw significant advancements with the development of specialized instruments. Edward Padgett and Earl Reese independently created mechanical dermatomes in the 1930s and 1940s, allowing for precise harvesting of split-thickness grafts. Innovations such as mesh grafting, introduced by Tanner and Vandeput in the mid-20th century, revolutionized the treatment of large burn injuries. Continuous refinements in wound bed preparation, infection control, and post-operative care have improved graft survival rates, solidifying skin grafting as a cornerstone of modern reconstructive surgery.

Skin grafting is indicated for a variety of conditions where significant skin loss or damage prevents primary closure and requires durable coverage. 

\begin{itemize}
  \item To cover deep partial-thickness or full-thickness burns that cannot heal spontaneously or with primary closure.
  \item For large traumatic wounds, such as degloving injuries or extensive lacerations, where significant skin loss has occurred.
  \item To close chronic non-healing ulcers, including diabetic foot ulcers and pressure ulcers, after meticulous debridement.
  \item Following the excision of large skin cancers when the resulting defect is too large for primary closure.
  \item For reconstruction after extensive debridement of severe infections like necrotizing fasciitis.
  \item To repair congenital defects requiring skin coverage, such as giant congenital nevi or omphalocele defects.
  \item To release and cover contractures, particularly those resulting from burns.
  \item To provide durable coverage over exposed vital structures like bone or tendon that cannot be covered by local tissue.
\end{itemize}

Skin grafting can serve as both a primary and a secondary procedure, depending on the clinical context. It is considered a primary procedure when immediate wound closure is necessary, such as in extensive full-thickness burns or acute traumatic injuries. In these cases, early excision of necrotic tissue followed by immediate grafting is often performed to achieve rapid wound closure and prevent complications.

More commonly, skin grafting functions as a secondary or salvage procedure after initial wound management, which may involve serial debridement and infection control. In these scenarios, the graft is applied to a clean, granulating wound bed that has been prepared. The decision to perform grafting as a primary or secondary procedure is guided by the urgency of wound closure, the condition of the wound bed, and the patient's overall health status.

\section*{Relevant Surgical Anatomy}
The skin, the largest organ of the body, consists of two primary layers: the epidermis and the dermis, which rests upon the subcutaneous fat. The epidermis serves as the outermost protective layer, primarily composed of keratinocytes, while the dermis contains collagen and elastin fibers, blood vessels, lymphatic vessels, nerve endings, and adnexal structures such as hair follicles, sebaceous glands, and sweat glands. The vascular supply to the skin is critical for graft viability, with the dermis containing both superficial and deep vascular plexuses that provide nourishment.

For split-thickness skin grafts (STSGs), only the epidermis and a portion of the superficial dermis are harvested, leaving behind residual dermal elements at the donor site that allow for spontaneous re-epithelialization. Common donor sites include the anterolateral thigh, buttocks, scalp, and abdomen, which are chosen for their large surface area and minimal scarring potential. Full-thickness skin grafts (FTSGs) involve harvesting the entire epidermis and dermis, typically from areas with skin laxity such as the groin, supraclavicular fossa, or postauricular region. Understanding the vascular supply, nerve distribution, and lymphatic drainage of these areas is crucial to minimize complications and ensure optimal graft take.

\begin{itemize}
  \item **Epidermis:** Outermost layer providing barrier function.
  \item **Dermis:** Contains blood vessels, nerves, and skin appendages.
  \item **Vascular Supply:** Superficial and deep plexuses critical for graft viability.
  \item **Donor Sites:** Common sites include thigh, abdomen, and scalp for STSGs; groin and supraclavicular region for FTSGs.
  \item **Nerve Relations:** Awareness of major nerves (e.g., lateral femoral cutaneous nerve) is essential to prevent injury.
\end{itemize}

\section*{Preoperative Workup \& Preparation}
Thorough pre-operative preparation is critical for optimizing the success of a skin graft. This involves a comprehensive assessment of the patient's overall health and meticulous optimization of the wound bed.

Key preparatory steps include:

\begin{itemize}
  \item \textbf{Comprehensive Patient Assessment:} Detailed history and physical examination to assess overall health and identify comorbidities.
  
  \item \textbf{Wound Bed Optimization:} Ensuring the recipient wound bed is clean, free of infection, and well-vascularized through debridement and infection control.
  
  \item \textbf{Laboratory and Imaging Studies:} Pre-operative blood tests and imaging studies as necessary to evaluate the wound and underlying conditions.
  
  \item \textbf{Anesthesia Clearance:} Cardiac and pulmonary clearance for patients undergoing general anesthesia.
  
  \item \textbf{Medication Management:} Adjusting anticoagulants and diabetic medications as needed.
  
  \item \textbf{Pre-operative Antibiotics:} Administering prophylactic antibiotics within 60 minutes prior to incision.
  
  \item \textbf{Informed Consent:} Discussing the procedure, risks, and expected outcomes with the patient.
  
  \item \textbf{NPO Status:} Instructing patients to fast prior to surgery.
\end{itemize}

Careful consideration of contraindications and precautions is essential to ensure patient safety and optimize graft success.

\begin{itemize}
  \item \textbf{Situations requiring treatment, optimization, or an alternative reconstruction:}
    \begin{itemize}
      \item \textbf{Recipient Bed Problems:} Uncontrolled infection, inadequate vascularity, exposed avascular structures, or active malignancy generally require treatment, further wound-bed preparation, flap coverage, or an oncologic plan before definitive grafting.
      \item \textbf{Systemic Risk:} Severe illness may require optimization or a different anesthetic and reconstructive approach.
    \end{itemize}
  
  \item \textbf{Relative Contraindications and Precautions:}
    \begin{itemize}
      \item \textbf{Poor Nutritional Status:} Malnutrition significantly impairs wound healing.
      \item \textbf{Uncontrolled Diabetes Mellitus:} Poorly controlled blood glucose levels increase the risk of complications.
      \item \textbf{Smoking:} Nicotine impairs microcirculation and increases graft failure risk.
      \item \textbf{Certain Medications:} High-dose corticosteroids and immunosuppressants can delay healing.
      \item \textbf{Patient Non-compliance:} Inability to adhere to post-operative care can jeopardize graft survival.
      \item \textbf{Recipient Site with Excessive Motion:} Grafts over highly mobile areas require careful immobilization.
      \item \textbf{Previous Radiation Therapy:} Irradiated tissue has compromised healing capacity.
    \end{itemize}
\end{itemize}

\section*{Surgical Technique \& Procedure}
The surgical technique for skin grafting involves several meticulous steps to ensure optimal graft take and patient outcomes. Anesthesia is typically general, but regional or local anesthesia may be used for smaller grafts. The patient is positioned to allow sterile access to both the recipient and donor sites.

The recipient wound bed is prepared by debriding necrotic tissue and achieving hemostasis. Donor site selection is crucial; for STSGs, common donor sites include the thigh and abdomen, while FTSGs are harvested from areas with skin laxity.

The key operative steps are as follows:

\begin{enumerate}
  \item \textbf{Recipient Site Preparation:} Thorough debridement of non-viable tissue and meticulous hemostasis.
  
  \item \textbf{Donor Site Preparation:} The donor site is prepped and draped, with a template used for FTSGs to mark the excision area.
  
  \item \textbf{Graft Harvesting:}
      \begin{itemize}
        \item For STSGs, a dermatome is used to harvest a uniform layer of skin, typically 0.010-0.018 inches thick. The graft may be meshed to allow for drainage.
        \item For FTSGs, the skin is excised with a scalpel, ensuring the entire thickness is included.
      \end{itemize}
  
  \item \textbf{Graft Application:} The harvested graft is transferred to the recipient bed, ensuring proper orientation and contact.
  
  \item \textbf{Graft Fixation:} The graft is secured using sutures or staples, and small fenestrations may be made to prevent fluid accumulation.
  
  \item \textbf{Dressing Application:} A non-adherent dressing is applied over the graft, followed by absorbent layers and a compressive dressing.
  
  \item \textbf{Donor Site Closure \& Dressing:} The donor site is dressed appropriately, with STSGs typically covered with an occlusive dressing.
  
  \item \textbf{Immobilization:} The grafted area is often immobilized to prevent disruption of the revascularization process.
\end{enumerate}

\section*{Complications \& Risks}
Skin grafting carries inherent risks, categorized into general surgical risks and procedure-specific risks.

\begin{itemize}
  \item \textbf{General Surgical Risks:}
    \begin{itemize}
      \item \textbf{Bleeding:} Hematoma formation can lead to graft failure.
      \item \textbf{Infection:} Both recipient and donor sites are susceptible to infection.
      \item \textbf{Anesthesia Complications:} Risks associated with anesthesia include allergic reactions and respiratory issues.
      \item \textbf{Pain:} Significant pain can occur at both sites, requiring effective management.
      \item \textbf{Scarring:} Scarring can be hypertrophic or keloid, leading to cosmetic concerns.
    \end{itemize}
  
  \item \textbf{Procedure-Specific Risks:}
    \begin{itemize}
      \item \textbf{Graft Failure:} Partial or complete non-take due to various factors.
      \item \textbf{Graft Contracture:} Can lead to functional impairment, especially over joints.
      \item \textbf{Donor Site Morbidity:} Pain, scarring, and delayed healing can occur at the donor site.
      \item \textbf{Pigmentary Changes:} Grafted skin may develop different pigmentation than surrounding skin.
      \item \textbf{Sensory Changes:} Altered sensation may occur in the grafted area.
      \item \textbf{Hair Growth:} Hair may continue to grow in grafted areas from hair-bearing donor sites.
      \item \textbf{Blistering or Fragility:} Thin grafts can be fragile and prone to blistering.
      \item \textbf{Recurrence of Underlying Condition:} If the primary cause of the wound is not addressed, recurrence may necessitate further intervention.
    \end{itemize}
\end{itemize}

\section*{Postoperative Care \& Recovery}
The post-operative period following a skin graft is crucial for ensuring graft take and optimizing recovery. 

In the immediate post-operative phase, patients are monitored in the Post-Anesthesia Care Unit (PACU) for recovery from anesthesia and pain control. The grafted area is protected with dressings and often immobilized to prevent movement. Pain management is prioritized, often involving a combination of oral and intravenous analgesics. 

Over the first 1-2 weeks, the viability of the graft is assessed. The first dressing change typically occurs between post-operative days 3-7, allowing evaluation of graft adherence. A successful graft will appear pink and adherent, while areas of failure may show signs of necrosis. Split-thickness donor sites usually heal within 1-3 weeks, while full-thickness donor sites require longer for complete healing. 

Long-term recovery involves scar management and physical therapy to maintain range of motion and prevent contractures. The grafted skin will initially be fragile and sensitive to sun exposure, requiring diligent protection. Sensory recovery can take months, and pigmentary changes are common.

During the skin graft procedure, the surgeon documents the characteristics of the recipient wound bed after debridement, including size, depth, and vascularity. At the donor site, the quality and quantity of skin harvested are noted. These findings are recorded in the operative report, including graft dimensions and any intraoperative complications.

Pathology reports are typically associated with the underlying lesion necessitating the graft rather than the graft itself. For example, if the graft follows skin cancer excision, the pathology report will confirm the type of malignancy and other prognostic features. In cases of chronic ulcers, pathology may reveal chronic inflammation or signs of infection. If a graft fails, the pathology of the failed tissue may show necrosis or inflammation, providing insight into the cause of failure.

A normal, successful postoperative course following a skin graft is characterized by complete integration and viability of the transferred skin. Within the first week, the graft should begin to revascularize, appearing pink and adherent. Pain should gradually decrease with appropriate analgesia. By 2-3 weeks, the donor site should be fully re-epithelialized, and the recipient site will be covered by the viable graft.

Over subsequent weeks to months, the grafted skin matures, becoming more durable, although it may retain differences in texture and color compared to surrounding skin. Patients should experience a progressive resolution of symptoms and a gradual return to improved function and cosmesis. The ultimate goal is complete wound coverage and restoration of skin barrier function, allowing the patient to resume daily activities with minimal impairment.

Postoperative care and follow-up are critical for ensuring graft take and managing complications.

\begin{itemize}
  \item \textbf{Regular Post-operative Visits:} Scheduled appointments to monitor graft take and assess healing progress.
  
  \item \textbf{Dressing Changes:} Regular changes of dressings for both graft and donor site, initially by medical staff.
  
  \item \textbf{Suture/Staple Removal:} Typically removed 7-14 days post-operatively, depending on healing progress.
  
  \item \textbf{Scar Management:} Initiation of scar massage and application of silicone sheets to minimize scarring.
  
  \item \textbf{Physical or Occupational Therapy:} Therapy may be prescribed to maintain range of motion and prevent contractures.
  
  \item \textbf{Sun Protection:} Strict avoidance of sun exposure and use of high-SPF sunscreen on grafted areas.
  
  \item \textbf{Warning Signs Education:} Patients are educated on signs of infection or graft loss to report immediately.
\end{itemize}

Adherence to these instructions is paramount for successful long-term outcomes.

\section*{Outcomes \& Prognosis}
Skin grafting is a widely performed reconstructive procedure, with its incidence directly correlated to the prevalence of conditions necessitating significant wound closure. Burns represent a major indication, with hundreds of thousands of burn injuries occurring annually in the United States, many requiring skin grafting. Traumatic injuries, including degloving injuries and large soft tissue defects, also contribute to the need for grafts. Chronic non-healing ulcers, such as diabetic foot ulcers and pressure ulcers, affect millions globally, with many requiring surgical intervention.

While skin grafts are performed across all age groups, the elderly population exhibits a higher incidence of chronic ulcers, whereas younger individuals are more frequently affected by traumatic injuries and burns. Mortality directly attributable to skin grafting is low, typically associated with severe underlying conditions or complications related to anesthesia. Morbidity primarily encompasses graft failure, donor site complications, and long-term issues such as contracture and scarring.

Graft take is most likely when the wound bed is clean, vascularized, stable, and appropriately immobilized, but failure can require further debridement or reconstruction. Outcomes vary with graft thickness, site, wound cause, vascularity, infection, smoking, nutrition, and comorbidity. Grafts restore coverage but may differ from surrounding skin in color, sensation, durability, and contour.

Recurrence of the wound requiring repeat grafting can occur, especially in patients with chronic ulcers if the underlying pathology is not adequately managed. Prognostic factors for successful graft take include a clean, well-vascularized recipient bed, absence of infection, meticulous surgical technique, and diligent post-operative care. Long-term prognosis involves managing potential complications such as hypertrophic scarring and pigmentary changes, which can impact both cosmesis and function.

\section*{References}
\begin{enumerate}
  \item MedlinePlus. Skin graft. U.S. National Library of Medicine; 2024. Available at: medlineplus.gov/skingraft.html. Accessed October 26, 2024.
  
  \item Townsend CM, Beauchamp RD, Evers BM, Mattox KL, editors. Sabiston Textbook of Surgery: The Biological Basis of Modern Surgical Practice. 22nd ed. Elsevier; 2022.
  
  \item American Society of Plastic Surgeons. Skin Grafting. Available at: plasticsurgery.org/reconstructive-procedures/skin-grafting. Accessed October 26, 2024.
  
  \item UpToDate. Basic principles of skin grafting. Available at: uptodate.com/contents/basic-principles-of-skin-grafting. Accessed October 26, 2024.
\end{enumerate}

\clearpage
